\documentclass[table]{beamer}
%[]中可以使用handout、trancompress等参数

%指定package
\usepackage{subfigure}
\usepackage{manfnt}%%% Dangerous Bend Symbols}\dbend \lhdbend \reversedvideodbend \textdbend \textlhdbend
\usepackage{multicol}
\usepackage{bbding}% 手势 \HandRight \HandLeft %\FiveStar \FourStar \SixStar
\usepackage{wasysym}
\usepackage{wasysym}
\usepackage{mflogo}

%指定beamer的模式与主题
\mode<presentation>
{
  \usetheme{CambridgeUS}
%Good themes
%\usetheme{CambridgeUS}
%\usetheme{Antibes}
%\usetheme{Madrid}
%\usetheme{Berkeley}
%\usetheme{Boadilla}

%\usecolortheme{default}
%\usecolortheme{orchid}
%\usecolortheme{whale}
\usecolortheme{spruce}
%\usefonttheme{professionalfonts}
}

%这里还可以选择别的主题：Bergen, Boadilla, Madrid, AnnArbor, CambridgeUS, Pittsburgh, Rochester, Warsaw, ...
%有导航栏的Antibes, JuanLesPins, Montpellier, ...
%有内容的Berkeley, PaloAlto, Goettingen, Marburg, Hannover, ...
%有最小导航栏的Berlin, Ilmenau, Dresden, Darmstadt, Frankfurt, Singapore, Szeged, ...
%有章和节表单的Copenhagen, Luebeck, Malmoe, Warsaw, ...

%\usecolortheme{default}
%设置内部颜色主题(这些主题一般改变block里的颜色);这个主题一般选择动物来命名
%这里还可以选择别的颜色主题,如默认的和有特别目的的颜色主题default,structure,sidebartab,全颜色主题albatross,beetle,crane,dove,fly,seagull,wolverine,beaver

%\usecolortheme{orchid}
%设置外部颜色主题(这些主题一般改变title里的颜色);这个主题一般选择植物来命名
%这里还可以选择别的颜色主题,如默认的和有特别目的的颜色主题lily,orchid,rose

%\usecolortheme{whale}
%设置字体主题;这个主题一般选择海洋动物来命名
%这里还可以选择别的颜色主题,如默认的和有特别目的的颜色主题whale,seahorse,dolphin

%\usefonttheme{professionalfonts}
%类似的还可以定义structurebold,structuresmallcapsserif,professionalfonts


% 控制 beamer 的风格,可以根据自己的爱好修改
%\usepackage{beamerthemesplit} %使用 split 风格
%\usepackage{beamerthemeshadow} %使用 shadow 风格
%\usepackage[width=2cm,dark,tab]{beamerthemesidebar}


% 设定英文字体
\usepackage{fontspec}
\usepackage{ listings}
\setmainfont{DejaVu Sans}
%\setmainfont{Times New Roman}
%\setsansfont{Arial}
%\setmonofont{Courier New}

% 设定中文字体
\usepackage[BoldFont,SlantFont,CJKchecksingle,CJKnumber]{xeCJK}
%\setCJKmainfont[BoldFont={Adobe Heiti Std},ItalicFont={Adobe Kaiti Std}]{Adobe Song Std}
\setCJKmainfont[BoldFont={WenQuanYi Micro Hei},ItalicFont={WenQuanYi Micro Hei}]{WenQuanYi Micro Hei}
\setCJKsansfont{WenQuanYi Micro Hei}
\setCJKmonofont{WenQuanYi Micro Hei}
\punctstyle{hangmobanjiao}

\defaultfontfeatures{Mapping=tex-text}
\usepackage{xunicode}
\usepackage{xltxtra}

\XeTeXlinebreaklocale "zh"
\XeTeXlinebreakskip = 0pt plus 1pt minus 0.1pt

\usepackage{setspace}
\usepackage{booktabs}
\usepackage{colortbl,xcolor}
\usepackage{hyperref}
%\hypersetup{xetex,bookmarksnumbered=true,bookmarksopen=true,pdfborder=1,breaklinks,colorlinks,linkcolor=blue,filecolor=black,urlcolor=cyan,citecolor=green}
\hypersetup{xetex,bookmarksnumbered=true,bookmarksopen=true,pdfborder=1,breaklinks,colorlinks,linkcolor=cyan,filecolor=black,urlcolor=blue,citecolor=green}

% 插入图片
\usepackage{graphicx}
% 指定存储图片的路径（当前目录下的figures文件夹）
\graphicspath{{figures/}}

% 可能用到的包
\usepackage{amsmath,amssymb}
\usepackage{multimedia}
\usepackage{multicol}

% 定义一些自选的模板，包括背景、图标、导航条和页脚等，修改要慎重
% 设置背景渐变由10%的红变成10%的结构颜色
%\beamertemplateshadingbackground{red!10}{structure!10}
%\beamertemplatesolidbackgroundcolor{white!90!blue}
% 使所有隐藏的文本完全透明、动态，而且动态的范围很小
\beamertemplatetransparentcovereddynamic
% 使itemize环境中变成小球，这是一种视觉效果
\beamertemplateballitem
% 为所有已编号的部分设置一个章节目录，并且编号显示成小球
\beamertemplatenumberedballsectiontoc
% 将每一页的要素的要素名设成加粗字体
\beamertemplateboldpartpage

% item逐步显示时,使已经出现的item、正在显示的item、将要出现的item呈现不同颜色
\def\hilite<#1>{
 \temporal<#1>{\color{gray}}{\color{blue}}
    {\color{blue!25}}
}

% 自定义彩色块状结构的颜色
\setbeamercolor{bluecyanbgcolor}{fg=blue,bg=cyan}
\setbeamercolor{redblackbgcolor}{fg=red,bg=black}
\setbeamercolor{greenwhitebgcolor}{fg=green,bg=white}
\setbeamercolor{blueblackbgcolor}{fg=blue,bg=black}
\setbeamercolor{yellowmagentabgcolor}{fg=yellow,bg=green}
\setbeamercolor{blackorangebgcolor}{fg=black,bg=orange}
\setbeamercolor{magentabgcolor}{fg=,bg=magenta}



% 在表格、图片等得标题中显示编号
\setbeamertemplate{caption}[numbered]

% 打开PDF后直接全屏
\hypersetup{pdfpagemode={FullScreen}}


% 使用 \part,\section,\subsection 等命令组织文档结构
% 使用 \frame 命令制作幻灯片

\begin{document}

\logo{\includegraphics[height=0.09\textwidth]{ihs.jpg}}
\title[Beamer模版演示]{使用基于\XeTeX 的Beamer制作幻灯片的模版}
\author[Yixf]{Yi Xianfu}
\institute[IHS, SIBS, CAS]{Institute of Health Sciences\\ Shanghai Institutes for Biological Sciences\\ Chinese Academy of Sciences}
\date{\today}

% 定义目录页
\AtBeginPart{
  \frame{
    \frametitle{\partpage}
    \begin{multicols}{2}
% 如果目录过长，可以打开这个选项分两栏显示
      \tableofcontents
% 使用这个命令自动生成目录
    \end{multicols}
  }  
}

% 在每个Section前都会加入的Frame
\AtBeginSection[]
{
  \begin{frame}<beamer>
    \frametitle{提纲}
	\setcounter{tocdepth}{2}
    \tableofcontents[currentsection,currentsubsection]
  \end{frame}
}
% 在每个Subsection前都会加入的Frame
\AtBeginSubsection[]
{
  \begin{frame}<beamer>
%\begin{frame}<handout:0>
% handout:0 表示只在手稿中出现
    \frametitle{提纲}
	\setcounter{tocdepth}{2}
    \tableofcontents[currentsection,currentsubsection]
% 显示在目录中加亮的当前章节
  \end{frame}
}

\begin{frame}
  \titlepage
\end{frame}

\begin{frame}[plain]
  \frametitle{提纲}
  \setcounter{tocdepth}{2}
  \tableofcontents
\end{frame}

\section{列表}
\subsection{有序列表}
\begin{frame}
	\frametitle{有序列表}
	  \begin{enumerate}[<+-|alert@+>]
		  \item 有序列表一
		  \item 有序列表二
		  \item 有序列表三
		  \item 有序列表四
		  \item 有序列表五
		  \item 有序列表六
		  \item 有序列表七
		  \item 有序列表八
		  \item 有序列表九
	  \end{enumerate}
\end{frame}
\subsection{无序列表}
\begin{frame}
	\begin{itemize}
		\hilite <1> \item 无序列表一
		\hilite <2-3> \item 无序列表二
		\hilite <3> \item 无序列表三
		\hilite <4-> \item 无序列表四
		\hilite <5> \item 无序列表五
		\hilite <-6> \item 无序列表六
		\hilite <7> \item 无序列表七
		\hilite <8> \item 无序列表八
		\hilite <9> \item 无序列表九
	\end{itemize}
\end{frame}

\section{图片}
\begin{frame}
	\frametitle{图片}
	\begin{figure}[htbp]
    \centering
    \includegraphics[width=8cm]{power.png}
    \caption{Powered by}
    \label{fig:power}
    \end{figure}
\end{frame}

\section{表格}
\begin{frame}
	\frametitle{表格}
    \begin{table}
    \centering
    \caption{一个表格示例}
    \rowcolors[]{1}{blue!20}{blue!10}
    \begin{tabular}{|c|c|c|c|}
    \hline
    \rowcolor{blue!50}学号&姓名&年龄&成绩\\
    \hline
    001 & 丁一 & 25 & 91\\
    002 & 刘二 & 24 & 89\\
    003 & 张三 & 25 & 90\\
    004 & 李四 & 24 & 95\\
    005 & 王五 & 25 & 93\\
    006 & 赵六 & 26 & 88\\
    007 & 钱七 & 25 & 86\\
    008 & 王八 & 24 & 89\\
    009 & 孙九 & 25 & 90\\
    \hline
    \end{tabular}
    \end{table}
\end{frame}

\section{块状结构}
\subsection{内置}
\begin{frame}
	\frametitle{block, definition, example}
	\begin{block}{块}
	这是一个block。
	\end{block}
	\pause
	\begin{definition}{定义：}
	这是一个definition。
	\end{definition}
	\pause
	\begin{example}{实例：}
	这是一个example。
	\end{example}
\end{frame}
\subsection{自定义}
\begin{frame}
	\begin{beamercolorbox}[rounded=true,shadow=true,wd=12cm]{bgcolor}
		这是一个自定义的彩色块状结构。
	\end{beamercolorbox}
\end{frame}

\section{分栏}
\begin{frame}
	\frametitle{分栏}
	\begin{columns}
	\column{3cm}
	这是第一栏的文字；栏宽3cm。
	\column{5cm}
	这是第二栏的文字；栏宽5cm。
	\end{columns}
\end{frame}

\section{公式}
\begin{frame}
	\frametitle{质能方程}
	\begin{equation}
	E=mc^2
	\label{emc}
	\end{equation}
\end{frame}

\begin{frame}
\end{frame}

\section*{Acknowledgements}
\begin{frame}
	\frametitle{Powered by}
	\begin{center}
		\includegraphics[width=9cm]{power.png}
	\end{center}
\end{frame}

\begin{frame}[plain]
	\begin{spacing}{1.5}
	\begin{center}
	\Huge{\textbf{Thanks for your attention!}}

	\Huge{\textit{Any questions?}}
\end{center}
\end{spacing}
\end{frame}

\end{document}
